\section{Open problems}
\label{p:open-problems}

Two problems are left open in the long form of this work, and both come
already equipped with a first body of machine-checked machinery.

\subsection{Spectral branch: global classification on the Brillouin torus}
Local infrared analysis says nothing about the \emph{global} structure
of the discrete quasi-energy. Two facts are verified: that the
quasi-energy is circular ($\varepsilon\in[-\pi,\pi]$ in units of
$\Delta t$), and that the two sectors $\omega=0$ and $\omega=\pi$ both
host band-crossing points. For the canonical 1D walk
$\Uop(k) = \cos(k)\bbI - \ii\sin(k)\sigma_{z}$ a direct check confirms
$\Uop(0)=\bbI$ and $\Uop(\pi)=-\bbI$, placing band touchings at both
$\omega=0$ and $\omega=\pi$. In three dimensions the eight-corner family
of the cubic Brillouin zone admits an assignment of $\chi=\pm 1$ to its
eight vertices with total sum zero and a $4$-and-$4$ distribution
between the two quasi-energy sectors --- a satisfiable SMT problem when
stated correctly, and unsatisfiable when one demands an imbalance such
as $5$ vs.\ $3$
(\texttt{validators/brillouin\_global.py}).

The open question is a \emph{complete} classification of the special
points of the spectrum of $\Uop(\vk)$ over the 3D Brillouin zone,
taking into account the exact symmetries of the discrete step. The
verification suite already covers the local algebra and the corner
combinatorics; what remains is the cohomological / global piece of the
picture, which we do not address here.

\subsection{Structural branch: protospace without reciprocal language}
The deeper open question is whether the underlying $\Proto$ admits a
formulation that does not assume a torus / reciprocal-space description
to begin with. Three concrete sub-tasks were left open in the long
form; each has a first verified answer.

\paragraph{(E1) Locality on non-periodic graphs.}
A tight-binding Hamiltonian on a chain of $n$ sites and on a
deliberately irregular graph of $6$ sites with mixed vertex degrees
yields $\Uop_{\Delta t} = \bbI - \ii\Delta t\,\Hop + O(\Delta t^{2})$
with the nearest-neighbour support pattern dictated by the graph
adjacency, without invoking Fourier analysis.

\paragraph{(E2) Chirality robust under local disorder.}
The discrete chiral symmetry $\Gamma$ anticommutes with a
bond-disordered SSH Hamiltonian
$(t_{1},t_{2},t_{1}{+}\varepsilon,t_{2})$ exactly, forcing the spectrum
to remain $\{-E,E\}$-symmetric without invoking translation invariance.

\paragraph{(E3) Index without torus.}
For a bipartite finite graph $G=(A\sqcup B,E)$, the nullity of the
adjacency is bounded below by $\lvert\lvert A\rvert-\lvert B\rvert\rvert$
and saturates on a maximum matching.

The conceptual significance of (E3) deserves emphasis. The conventional
Nielsen--Ninomiya argument runs over the Brillouin torus
$T^{d}$~\cite{NielsenNinomiya1981a,NielsenNinomiya1981b}: it counts band
crossings homotopically as obstructions on $T^{d}$, and so depends
essentially on the reciprocal-space description. The bipartite index
bound is the same chirality balance for a finite graph $G$ with no
group of translations and no torus to homotope over; the argument
reduces to linear algebra on the adjacency. This is the algebraic
backbone needed if one ever wants to phrase \enquote{protospace} as a
piece of structure broader than a periodic lattice --- for instance
along the causal-set
line~\cite{BombelliLeeMeyerSorkin1987,Sorkin2003} --- without smuggling
the torus back in.

All three sub-tasks have machine-checked tests
(\texttt{validators/rama\_estructural.py}), so the structural branch is
genuinely \emph{open but well-posed}: the algebra and combinatorics
necessary to formulate it on arbitrary graphs is in place; what is
missing is a principle (categorial, topological, or otherwise) that
selects \emph{which} graphs are the ones a protospace can look like. A
companion paper devoted entirely to the structural branch is in
preparation.
